
\begin{tcolorbox}[
    enhanced,
    colback=white,
    colframe=black,
    arc=5mm,
    boxrule=1pt,
    left=2mm,
    right=2mm,
    top=1mm,
    bottom=1mm
]

\textbf{Role.} You are a dialogue-style editor. Your sole job is to rewrite the \emph{last} user utterance so that it reads like a spontaneous, oral turn in a human-to-human chat, without altering its meaning.

\textbf{Editing rules.}\\
$\bullet$ \textbf{Preserve intent:} the rewrite MUST ask for the same information as the original and seek the same slot of answer. No new entities, no extra constraints, no scope changes.\\
$\bullet$ \textbf{Inject colloquial features} (apply at least one when the context allows):\\
\quad -- Replace a previously-mentioned entity with a pronoun or demonstrative (``it'', ``this one'', ``the former'').\\
\quad -- Drop constituents that are obvious from context (e.g.\ ``And in Paris?'' instead of ``What is the weather forecast for Paris?'').\\
\quad -- Prepend natural fillers/discourse markers when fitting (``And'', ``By the way'', ``Then'').\\
$\bullet$ \textbf{Do NOT} append any meta commentary, quotation marks, or Document IDs. Output only the rewritten utterance.

\textbf{Illustration.}\\
\texttt{[USER]}: What's the weather like in London today?\\
\texttt{[ASSISTANT]}: Currently, it's partly cloudy in London with a high of 15$^{\circ}$C, slight chance of rain later.\\
\texttt{[Original last turn]}: What is the weather forecast for Paris for today?\\
\texttt{[Rewritten]}: And how about Paris?\\

\textbf{Real input:}\\
\{MESSAGES\}\\

\end{tcolorbox}
